%%%%%%%%%%%%%%%%%%%%%%%%%%%%%%%%%%%%%%%%%%%%%%%%%%%%%%%%%%%%%%%%%%%%%%%%
%%% LaTeX Report for CS6002 Project (AAMAS format)
%%%%%%%%%%%%%%%%%%%%%%%%%%%%%%%%%%%%%%%%%%%%%%%%%%%%%%%%%%%%%%%%%%%%%%%%

\documentclass[sigconf]{aamas}

\usepackage{balance}
\usepackage{amsmath}
\usepackage{amssymb}
\usepackage{amsfonts}
\usepackage{mathtools}
\usepackage{algorithm}
\usepackage{algorithmic}
\usepackage{bm}

%%%%%%%%%%%%%%%%%%%%%%%%%%%%%%%%%%%%%%%%%%%%%%%%%%%%%%%%%%%%%%%%%%%%%%%%
%%% Copyright block
%%%%%%%%%%%%%%%%%%%%%%%%%%%%%%%%%%%%%%%%%%%%%%%%%%%%%%%%%%%%%%%%%%%%%%%%

\setcopyright{ifaamas}
\acmConference[AAMAS '25]{Proc.\@ of the 24th International Conference
on Autonomous Agents and Multiagent Systems (AAMAS 2025)}{May 19 -- 23, 2025}
{Detroit, Michigan, USA}{}
\copyrightyear{2025}
\acmYear{2025}
\acmDOI{}
\acmPrice{}
\acmISBN{}

\acmSubmissionID{}

%%%%%%%%%%%%%%%%%%%%%%%%%%%%%%%%%%%%%%%%%%%%%%%%%%%%%%%%%%%%%%%%%%%%%%%%
%%% Title and Authors
%%%%%%%%%%%%%%%%%%%%%%%%%%%%%%%%%%%%%%%%%%%%%%%%%%%%%%%%%%%%%%%%%%%%%%%%

\title{Nash Equilibria in SEAD via Counterfactual Regret Minimisation Under Unknown Attacker Types}

%%% TODO: Replace with actual author names and affiliations
\author{Author 1}
\affiliation{
  \institution{IIT Madras}
  \city{Chennai}
  \country{India}}
\email{author1@iitm.ac.in}

\author{Author 2}
\affiliation{
  \institution{IIT Madras}
  \city{Chennai}
  \country{India}}
\email{author2@iitm.ac.in}

%%%%%%%%%%%%%%%%%%%%%%%%%%%%%%%%%%%%%%%%%%%%%%%%%%%%%%%%%%%%%%%%%%%%%%%%
%%% Abstract
%%%%%%%%%%%%%%%%%%%%%%%%%%%%%%%%%%%%%%%%%%%%%%%%%%%%%%%%%%%%%%%%%%%%%%%%

\begin{abstract}
Suppression of Enemy Air Defences (SEAD) involves a repeated, simultaneous-move engagement between an attacking force and a defending surface-to-air missile (SAM) battery under resource depletion and imperfect information. Classical Stackelberg Security Game models are structurally inappropriate for SEAD because neither side can credibly commit to a strategy. We model the SEAD engagement as a Normal-Form Game with Sequential Strategies (NFGSS) and compute Nash equilibria using CFR+. We then address the operationally realistic setting where the defender does not know the set of attacker types \emph{a priori}: the true number of types is bounded by a known upper bound $K_{\max}$, and the defender must learn the attacker population online. Adapting the online learning framework of Balcan et al.\ to our simultaneous-move setting, we design an explore--exploit algorithm that maintains a growing catalogue of observed attacker types and periodically re-solves the induced game via CFR+. We aim to show empirically that the resulting defender strategy achieves no-regret behaviour, with cumulative regret sublinear in the number of engagement episodes.
\end{abstract}

\keywords{Security Games, SEAD, Counterfactual Regret Minimisation, Online Learning, Nash Equilibrium, Unknown Attacker Types}

%%%%%%%%%%%%%%%%%%%%%%%%%%%%%%%%%%%%%%%%%%%%%%%%%%%%%%%%%%%%%%%%%%%%%%%%
%%% Custom commands
%%%%%%%%%%%%%%%%%%%%%%%%%%%%%%%%%%%%%%%%%%%%%%%%%%%%%%%%%%%%%%%%%%%%%%%%

\newcommand{\Satt}{S^{\mathrm{att}}}
\newcommand{\Sdef}{S^{\mathrm{def}}}
\newcommand{\Aatt}{\mathcal{A}^{\mathrm{att}}}
\newcommand{\Adef}{\mathcal{A}^{\mathrm{def}}}
\newcommand{\pdeceive}{p_{\mathrm{dec}}}
\newcommand{\phit}{p_{\mathrm{hit}}}
\newcommand{\Kmax}{K_{\max}}
\newcommand{\E}{\mathbb{E}}
\newcommand{\R}{\mathbb{R}}
\newcommand{\indicator}{\mathbb{I}}

%%%%%%%%%%%%%%%%%%%%%%%%%%%%%%%%%%%%%%%%%%%%%%%%%%%%%%%%%%%%%%%%%%%%%%%%

\begin{document}

\pagestyle{fancy}
\fancyhead{}

\maketitle

%%%%%%%%%%%%%%%%%%%%%%%%%%%%%%%%%%%%%%%%%%%%%%%%%%%%%%%%%%%%%%%%%%%%%%%%
\section{Introduction}
%%%%%%%%%%%%%%%%%%%%%%%%%%%%%%%%%%%%%%%%%%%%%%%%%%%%%%%%%%%%%%%%%%%%%%%%

Suppression of Enemy Air Defences (SEAD) is among the most tactically critical engagements in modern air warfare. An attacking force---comprising strike aircraft, electronic warfare platforms, and expendable decoys---must neutralise or suppress a defender's integrated air defence system (IADS) before the main strike package can proceed. Both sides operate under \emph{severe resource constraints} (the attacker has a finite decoy budget; the defender has a finite missile inventory), \emph{imperfect information} (the defender cannot distinguish real aircraft from decoys with certainty), and \emph{simultaneous decision-making} over multiple rounds as resources deplete.

The dominant paradigm for modelling security resource allocation is the \emph{Stackelberg Security Game} (SSG)~\cite{Tambe2012}, in which a defender commits to a mixed strategy and an attacker best-responds with full knowledge of that commitment. This leader--follower structure is justified in civilian security (e.g., patrol schedules leak over time), but breaks down in SEAD: a SAM battery commander does not publish an engagement doctrine, both sides choose actions simultaneously in each round, and the strategy space itself evolves as resources deplete.

Several papers address the problem of \emph{learning} the attacker's type through repeated interaction~\cite{Letchford2009, Marecki2012, Blum2014}, but these are restricted to the Stackelberg commitment model with a single or known set of attacker types. Balcan et al.~\cite{Balcan2015} introduced the first no-regret online learning framework for repeated Stackelberg security games with multiple, adversarially chosen attacker types from a known set~$K$, achieving regret bounds of $O(\sqrt{T n^2 k \log nk})$ under full information and $O(T^{2/3} nk \log^{1/3}(nk))$ under partial information. Meanwhile, Zinkevich et al.~\cite{Zinkevich2007} introduced counterfactual regret minimisation (CFR) for computing Nash equilibria in extensive-form games with imperfect information, and Lis\'{y} et al.~\cite{Lisy2016} adapted CFR to sequential security games via the NFGSS framework.

Our work synthesises these threads. We model SEAD as an NFGSS and solve for Nash equilibria using CFR+, then extend the framework to the operationally realistic case where the defender does not know the attacker's type set \emph{a priori} but has an upper bound~$\Kmax$ on its size. We design a hybrid explore--exploit algorithm inspired by Balcan et al.~\cite{Balcan2015}, adapted from the Stackelberg setting to our simultaneous-move CFR-based framework, and evaluate it empirically.

%%%%%%%%%%%%%%%%%%%%%%%%%%%%%%%%%%%%%%%%%%%%%%%%%%%%%%%%%%%%%%%%%%%%%%%%
\section{Preliminaries}
\label{sec:prelim}
%%%%%%%%%%%%%%%%%%%%%%%%%%%%%%%%%%%%%%%%%%%%%%%%%%%%%%%%%%%%%%%%%%%%%%%%

We first establish the formal framework for extensive-form games, regret minimisation, and Stackelberg security games that underpins our model. Our notation follows Zinkevich et al.~\cite{Zinkevich2007} for the game-theoretic foundations and Balcan et al.~\cite{Balcan2015} for the online learning extension.

%-----------------------------------------------------------------------
\subsection{Extensive-Form Games with Imperfect Information}
\label{sec:efg}
%-----------------------------------------------------------------------

Following~\cite{Osborne1994, Zinkevich2007}, a finite extensive-form game with imperfect information consists of:
\begin{itemize}
    \item A finite set $N = \{1, 2\}$ of \textbf{players} (in our setting, the attacker and the defender), plus a chance player~$c$.
    \item A finite set $H$ of \textbf{histories} (sequences of actions). The empty sequence is in $H$, and every prefix of a sequence in $H$ is also in $H$. The set of \textbf{terminal histories} is $Z \subseteq H$. For any non-terminal history $h \in H \setminus Z$, the set of available actions is $A(h) = \{a : (h, a) \in H\}$.
    \item A \textbf{player function} $P : H \setminus Z \to N \cup \{c\}$ assigning each non-terminal history to the player who acts.
    \item For each player $i \in N$, a partition $\mathcal{I}_i$ of $\{h \in H : P(h) = i\}$ into \textbf{information sets}, with $A(h) = A(h')$ whenever $h, h'$ belong to the same $I_i \in \mathcal{I}_i$. We write $A(I_i)$ for this common action set.
    \item A \textbf{utility function} $u_i : Z \to \R$ for each player $i \in N$. The game is \textbf{zero-sum} if $u_1(z) = -u_2(z)$ for all $z \in Z$.
\end{itemize}

A \textbf{strategy} $\sigma_i$ for player $i$ assigns a probability distribution over $A(I_i)$ to each information set $I_i \in \mathcal{I}_i$. A \textbf{strategy profile} is $\sigma = (\sigma_1, \sigma_2)$, with $\sigma_{-i}$ denoting all strategies except player~$i$'s. Let $\pi^{\sigma}(h)$ denote the probability of history $h$ occurring under $\sigma$, decomposed as $\pi^{\sigma}(h) = \pi_i^{\sigma}(h) \cdot \pi_{-i}^{\sigma}(h)$, where $\pi_i^{\sigma}(h)$ is player $i$'s contribution and $\pi_{-i}^{\sigma}(h)$ is the product of all other players' contributions (including chance). The expected utility of player $i$ under profile $\sigma$ is:
\begin{equation}
    u_i(\sigma) = \sum_{h \in Z} u_i(h) \, \pi^{\sigma}(h).
    \label{eq:expected_utility}
\end{equation}

%-----------------------------------------------------------------------
\subsection{Nash Equilibrium}
\label{sec:nash}
%-----------------------------------------------------------------------

A \textbf{Nash equilibrium} is a strategy profile $\sigma^*$ such that no player can improve their expected utility by unilateral deviation:
\begin{equation}
    u_i(\sigma^*) \geq \max_{\sigma_i' \in \Sigma_i} \, u_i(\sigma_i', \sigma_{-i}^*), \quad \forall \, i \in N.
    \label{eq:nash}
\end{equation}
An \textbf{$\epsilon$-Nash equilibrium} relaxes this to:
\begin{equation}
    u_i(\sigma^*) + \epsilon \geq \max_{\sigma_i' \in \Sigma_i} \, u_i(\sigma_i', \sigma_{-i}^*), \quad \forall \, i \in N.
    \label{eq:eps_nash}
\end{equation}

%-----------------------------------------------------------------------
\subsection{Regret Minimisation and CFR}
\label{sec:cfr}
%-----------------------------------------------------------------------

\textbf{Average overall regret.}\quad
Let $\sigma_i^t$ be the strategy used by player $i$ at iteration $t$. The average overall regret of player $i$ after $T$ iterations is~\cite{Zinkevich2007}:
\begin{equation}
    R_i^T = \frac{1}{T} \max_{\sigma_i^* \in \Sigma_i} \sum_{t=1}^{T} \Big( u_i(\sigma_i^*, \sigma_{-i}^t) - u_i(\sigma^t) \Big).
    \label{eq:overall_regret}
\end{equation}
By Theorem~2 of \cite{Zinkevich2007}, if both players' average overall regret is less than $\epsilon$, then the average strategy profile $\bar{\sigma}^T$ is a $2\epsilon$-Nash equilibrium.

\textbf{Counterfactual regret.}\quad
Define the \textbf{counterfactual value} of player $i$ at information set $I$ under profile $\sigma$ as~\cite{Zinkevich2007}:
\begin{equation}
    v_i(\sigma, I) = \frac{\sum_{h \in I, h' \in Z} \pi_{-i}^{\sigma}(h) \, \pi^{\sigma}(h, h') \, u_i(h')}{\pi_{-i}^{\sigma}(I)},
    \label{eq:cfv}
\end{equation}
where $\pi^{\sigma}(h, h')$ is the probability of transitioning from $h$ to $h'$ under $\sigma$, and $\pi_{-i}^{\sigma}(I) = \sum_{h \in I} \pi_{-i}^{\sigma}(h)$.

For each action $a \in A(I)$, let $\sigma|_{I \to a}$ denote the strategy identical to $\sigma$ except that action $a$ is always chosen at $I$. The \textbf{immediate counterfactual regret} for action $a$ at information set $I$ is:
\begin{equation}
    R_{i}^{T}(I, a) = \frac{1}{T} \sum_{t=1}^{T} \pi_{-i}^{\sigma^t}(I) \Big( u_i(\sigma^t|_{I \to a}, I) - u_i(\sigma^t, I) \Big).
    \label{eq:cfr_action}
\end{equation}
Let $R_{i}^{T,+}(I, a) = \max(R_i^T(I, a), 0)$. The key result of~\cite{Zinkevich2007} (Theorem~3) is that overall regret is bounded by the sum of positive counterfactual regrets:
\begin{equation}
    R_i^T \leq \sum_{I \in \mathcal{I}_i} R_{i,\mathrm{imm}}^{T,+}(I).
    \label{eq:cfr_bound}
\end{equation}

\textbf{CFR algorithm.}\quad
The strategy at iteration $T+1$ is set via regret matching~\cite{Zinkevich2007}:
\begin{equation}
    \sigma_i^{T+1}(I)(a) = \begin{cases}
        \displaystyle\frac{R_i^{T,+}(I, a)}{\sum_{a' \in A(I)} R_i^{T,+}(I, a')} & \text{if } \sum_{a'} R_i^{T,+}(I, a') > 0, \\[8pt]
        \displaystyle\frac{1}{|A(I)|} & \text{otherwise.}
    \end{cases}
    \label{eq:regret_matching}
\end{equation}
Under this rule, $R_{i,\mathrm{imm}}^{T}(I) \leq \Delta_{u,i} \sqrt{|A_i|} / \sqrt{T}$ for each information set, and therefore the average strategy profile converges to a Nash equilibrium at rate $O(1/\sqrt{T})$.

\textbf{CFR+.}\quad
CFR+ \cite{Tammelin2015} modifies CFR by clipping cumulative regrets to be non-negative and using linear averaging (iteration $t$ receives weight $t$), yielding dramatically faster empirical convergence while preserving the $O(1/\sqrt{T})$ theoretical guarantee.

%-----------------------------------------------------------------------
\subsection{Stackelberg Security Games}
\label{sec:ssg}
%-----------------------------------------------------------------------

For comparison, we recall the Stackelberg security game framework~\cite{Tambe2012, Balcan2015}. A \textbf{repeated Stackelberg security game} consists of:
\begin{itemize}
    \item A time horizon $T$ and a set of \textbf{targets} $\mathcal{N} = \{1, \ldots, n\}$.
    \item A defender who commits to a \textbf{coverage probability} vector $\mathbf{p} \in [0,1]^n$, where $p_i$ is the probability that target $i$ is defended.
    \item A set of attacker types $K = \{\alpha_1, \ldots, \alpha_k\}$, selected adversarially from a known set.
    \item \textbf{Utilities}: For defender, $U_d(i, \mathbf{p}) = u_d^c(i) p_i + u_d^u(i)(1 - p_i)$; for attacker type $\alpha_j$, $U_{\alpha_j}(i, \mathbf{p}) = u_{\alpha_j}^c(i) p_i + u_{\alpha_j}^u(i)(1 - p_i)$.
    \item Attacker $\alpha_j$ \textbf{best-responds}: $b_{\alpha_j}(\mathbf{p}) = \arg\max_{i \in \mathcal{N}} U_{\alpha_j}(i, \mathbf{p})$.
\end{itemize}

The \textbf{best mixed strategy in hindsight} against a sequence of attackers $\mathbf{a} = a_1, \ldots, a_T$ is~\cite{Balcan2015}:
\begin{equation}
    \mathbf{p}^* = \arg\max_{\mathbf{p}} \sum_{t=1}^{T} U_d(b_{a_t}(\mathbf{p}), \mathbf{p}).
    \label{eq:best_hindsight}
\end{equation}
The \textbf{regret} of an online algorithm that plays $\mathbf{p}_t$ at step $t$ is:
\begin{equation}
    \mathrm{Regret}(T) = \sum_{t=1}^{T} U_d(b_{a_t}(\mathbf{p}^*), \mathbf{p}^*) - \E\!\left[\sum_{t=1}^{T} U_d(b_{a_t}(\mathbf{p}_t), \mathbf{p}_t)\right].
    \label{eq:regret_ssg}
\end{equation}
An algorithm is \textbf{no-regret} if $\mathrm{Regret}(T) / T \to 0$ as $T \to \infty$.

%%%%%%%%%%%%%%%%%%%%%%%%%%%%%%%%%%%%%%%%%%%%%%%%%%%%%%%%%%%%%%%%%%%%%%%%
\section{Problem Formulation}
\label{sec:problem}
%%%%%%%%%%%%%%%%%%%%%%%%%%%%%%%%%%%%%%%%%%%%%%%%%%%%%%%%%%%%%%%%%%%%%%%%

We now formally define the SEAD engagement model and the two problems we address: (i)~computing Nash equilibria in the base SEAD game via CFR+, and (ii)~online learning of defender strategies when the attacker type set is unknown.

%-----------------------------------------------------------------------
\subsection{The SEAD Game Model}
\label{sec:sead_model}
%-----------------------------------------------------------------------

We model the SEAD engagement as a two-player zero-sum \textbf{Normal-Form Game with Sequential Strategies} (NFGSS)~\cite{Lisy2016}. Each player's sequential decision process is captured by a finite acyclic Markov Decision Process (MDP), and the combination of both MDPs defines an extensive-form game with imperfect information.

\subsubsection{Attacker MDP}
The attacker's MDP is the tuple $\langle \Satt, \Aatt, T^{\mathrm{att}}, R_{\mathrm{att}} \rangle$:
\begin{itemize}
    \item \textbf{States} $\Satt$: Each state $s^{\mathrm{att}}$ is a tuple $(\tau, d, r, \mathbf{o}^{\mathrm{att}})$, where $\tau \in \{1, \ldots, R_{\max}\}$ is the current round, $d \in \{0, \ldots, D\}$ is the number of decoys remaining, $r \in \{0, \ldots, A_{\mathrm{real}}\}$ is the number of real aircraft remaining, and $\mathbf{o}^{\mathrm{att}}$ encodes the attacker's observation of previous defender responses. Here $D$ denotes the initial decoy budget and $A_{\mathrm{real}}$ the initial number of real aircraft.
    \item \textbf{Actions} $\Aatt(s)$: At state $s$, the attacker chooses from:
    \begin{itemize}
        \item \textsc{Deploy-Decoy}$(j)$: deploy a decoy to sector $j \in \{1, \ldots, J\}$ (requires $d > 0$),
        \item \textsc{Deploy-Real}$(j)$: deploy a real aircraft to sector $j$ (requires $r > 0$),
        \item \textsc{Hold}: take no action this round,
        \item \textsc{Ceasefire}: signal willingness to disengage.
    \end{itemize}
    The set of available actions depends on the resource counts $(d, r)$.
    \item \textbf{Transitions} $T^{\mathrm{att}}$: Deploying a decoy reduces $d$ by $1$ deterministically. The decoy is detected as fake with probability $\pdeceive$ (in which case the defender does not expend a missile) and is mistaken for a real aircraft with probability $1 - \pdeceive$. Deploying a real aircraft leads to engagement: the aircraft survives with probability $1 - \phit$ (which depends on whether the defender fires and the defender's radar mode).
\end{itemize}

\subsubsection{Defender MDP}
The defender's MDP is the tuple $\langle \Sdef, \Adef, T^{\mathrm{def}}, R_{\mathrm{def}} \rangle$:
\begin{itemize}
    \item \textbf{States} $\Sdef$: Each state $s^{\mathrm{def}}$ is a tuple $(\tau, m, \mathbf{r}_{\mathrm{tracks}}, \mathbf{o}^{\mathrm{def}})$, where $\tau$ is the round, $m \in \{0, \ldots, M\}$ is the number of missiles remaining (with $M$ the initial missile inventory), $\mathbf{r}_{\mathrm{tracks}}$ describes the currently active radar tracks, and $\mathbf{o}^{\mathrm{def}}$ encodes observations of attacker actions.
    \item \textbf{Actions} $\Adef(s)$: At state $s$, the defender chooses from:
    \begin{itemize}
        \item \textsc{Fire}$(j)$: fire a missile at radar track $j$ (requires $m > 0$),
        \item \textsc{Track-Only}$(j)$: maintain radar lock on track $j$ without firing,
        \item \textsc{Conserve}: power down radar to avoid anti-radiation missiles,
        \item \textsc{Ceasefire}: accept disengagement.
    \end{itemize}
    \item \textbf{Transitions} $T^{\mathrm{def}}$: Firing reduces $m$ by $1$ deterministically. The engagement outcome (hit or miss) is stochastic, with hit probability depending on whether the target is a real aircraft or a decoy, and on the defender's radar mode.
\end{itemize}

\subsubsection{Joint Game and Payoffs}
At each round $\tau$, both players simultaneously select actions from their respective MDPs. The actions are revealed, engagement outcomes are resolved stochastically, and both players transition to their next states. The game terminates when $\tau = R_{\max}$, when both players signal ceasefire, or when either side's critical resources are fully depleted.

Let $z = (s_1^{\mathrm{att}}, a_1^{\mathrm{att}}, s_1^{\mathrm{def}}, a_1^{\mathrm{def}}, \ldots, s_{R_{\max}}^{\mathrm{att}}, a_{R_{\max}}^{\mathrm{att}}, s_{R_{\max}}^{\mathrm{def}}, a_{R_{\max}}^{\mathrm{def}})$ denote a terminal history. The terminal utility for the defender is:
\begin{equation}
    u_{\mathrm{def}}(z) = \sum_{\tau=1}^{R_{\max}} U\!\left(s_\tau^{\mathrm{att}}, a_\tau^{\mathrm{att}}, s_\tau^{\mathrm{def}}, a_\tau^{\mathrm{def}}\right),
    \label{eq:terminal_utility}
\end{equation}
where $U(\cdot)$ is the per-round marginal utility capturing the resource trade-offs of each engagement outcome (e.g., defender fires at decoy: defender loses a missile, attacker loses a low-cost decoy; defender fires at real aircraft and hits: attacker loses a high-value asset). The game is zero-sum: $u_{\mathrm{att}}(z) = -u_{\mathrm{def}}(z)$.

\subsubsection{NFGSS Structure}
Following Lis\'{y} et al.~\cite{Lisy2016}, this SEAD game is an NFGSS: each player's strategy space is defined by its MDP, and the simultaneous-move interaction at each round creates imperfect information (neither player observes the other's action before choosing their own within a round). The NFGSS converts to a chance-relaxed skew well-formed (CRSWF) extensive-form game with zero error terms, guaranteeing that CFR+ converges to a Nash equilibrium.

%-----------------------------------------------------------------------
\subsection{Problem 1: Nash Equilibrium Computation}
\label{sec:problem1}
%-----------------------------------------------------------------------

\textbf{Given} the SEAD NFGSS $\Gamma$ with parameters $(D, A_{\mathrm{real}}, M, J, R_{\max}, \pdeceive, \phit)$ and a \emph{known} set of attacker types $K = \{\alpha_1, \ldots, \alpha_k\}$ (each type corresponding to a distinct parameterisation of the attacker MDP), \textbf{compute} an $\epsilon$-Nash equilibrium $\sigma^* = (\sigma_{\mathrm{att}}^*, \sigma_{\mathrm{def}}^*)$ such that:
\begin{equation}
    u_i(\sigma^*) + \epsilon \geq \max_{\sigma_i' \in \Sigma_i} u_i(\sigma_i', \sigma_{-i}^*), \quad \forall \, i \in \{\mathrm{att}, \mathrm{def}\}.
    \label{eq:problem1}
\end{equation}

We solve this using CFR+ (Section~\ref{sec:cfr}). After $T_{\mathrm{iter}}$ iterations, the average strategy profile $\bar{\sigma}^{T_{\mathrm{iter}}}$ is a $2\epsilon$-Nash equilibrium with $\epsilon = O(1/\sqrt{T_{\mathrm{iter}}})$, where the constant depends on $|\mathcal{I}_i|$ and $|A_i|$ (the number of information sets and maximum action count).

%-----------------------------------------------------------------------
\subsection{Problem 2: Online Learning Under Unknown Attacker Types}
\label{sec:problem2}
%-----------------------------------------------------------------------

The base model (Problem~1) assumes the defender knows $K$. In operational SEAD, this assumption fails: the defender may face novel attack profiles---previously unseen decoy types, unexpected electronic warfare techniques, or unfamiliar force compositions. We now formalise this uncertainty.

\subsubsection{Repeated Game Setup}
Consider a repeated SEAD engagement over $T$ \textbf{episodes} (indexed by $t = 1, \ldots, T$). At each episode $t$:
\begin{enumerate}
    \item The defender selects a strategy $\sigma_{\mathrm{def}}^t$ (possibly randomised).
    \item An adversary selects an attacker type $\alpha_t$ from a \emph{true} (but unknown) set of types $K^* \subseteq \{\alpha_1, \ldots, \alpha_{\Kmax}\}$, where $|K^*| \leq \Kmax$ and $\Kmax$ is a known upper bound.
    \item The episode plays out: the attacker of type $\alpha_t$ plays the SEAD game against the defender's strategy $\sigma_{\mathrm{def}}^t$. The defender receives payoff $u_{\mathrm{def}}(\sigma_{\mathrm{def}}^t, \beta_{\alpha_t}(\sigma_{\mathrm{def}}^t))$, where $\beta_{\alpha_t}(\cdot)$ denotes the attacker's best response (or Nash response in the simultaneous-move setting).
    \item The defender receives \textbf{feedback} about the episode.
\end{enumerate}

We consider two feedback models, following~\cite{Balcan2015}:
\begin{itemize}
    \item \textbf{Full information}: After episode $t$, the defender observes the attacker type $\alpha_t$ (e.g., the exact decoy model used). Formally, the defender observes $\alpha_t$ and can compute $\beta_{\alpha_j}(\sigma)$ for any $\alpha_j$ and any $\sigma$.
    \item \textbf{Partial information}: The defender only observes the \emph{outcome} of the engagement (which targets were attacked, hits/misses), not the attacker's type directly. Formally, the defender observes only the terminal history $z_t$ from which $\alpha_t$ must be inferred.
\end{itemize}

\subsubsection{Catalogue-Based Learning}
The defender maintains a \textbf{growing catalogue} $C_t \subseteq \{\alpha_1, \ldots, \alpha_{\Kmax}\}$ of attacker types observed or inferred by episode~$t$. Initially $C_0 = \emptyset$ (or a small prior set). As new types are encountered, they are added to the catalogue. We define:
\begin{itemize}
    \item The \textbf{induced game} $\Gamma(C_t)$: the SEAD NFGSS restricted to attacker types in the current catalogue $C_t$. As $C_t$ grows toward $K^*$, the game $\Gamma(C_t) \to \Gamma(K^*)$.
    \item The \textbf{induced Nash equilibrium} $\sigma^*(C_t)$: the Nash equilibrium of $\Gamma(C_t)$, computed via CFR+.
\end{itemize}

\subsubsection{Explore--Exploit Strategy}
The defender follows an explore--exploit strategy parameterised by an exploration probability $\varepsilon_t$:
\begin{equation}
    \sigma_{\mathrm{def}}^t = \begin{cases}
        \sigma_{\mathrm{explore}} & \text{with probability } \varepsilon_t, \\
        \sigma^*(C_{t-1}) & \text{with probability } 1 - \varepsilon_t,
    \end{cases}
    \label{eq:explore_exploit}
\end{equation}
where $\sigma_{\mathrm{explore}}$ is a randomised strategy designed to elicit type-revealing responses from the attacker, and $\varepsilon_t$ is a decaying exploration probability (e.g., $\varepsilon_t = 1/\sqrt{t}$) that balances information acquisition against exploitation of the current equilibrium. Whenever $|C_t| > |C_{t-1}|$ (a new type is discovered), the defender re-solves the expanded game $\Gamma(C_t)$ using CFR+, warm-starting from the previous solution.

\subsubsection{Regret Definition}
We define regret relative to an oracle defender who knows $K^*$ from the start. Let $\sigma^*(K^*)$ denote the Nash equilibrium of the true game $\Gamma(K^*)$, and let $u_{\mathrm{def}}^*$ be the defender's expected per-episode payoff under $\sigma^*(K^*)$. The \textbf{cumulative regret} of our online algorithm after $T$ episodes is:
\begin{equation}
    \mathrm{Regret}(T) = \sum_{t=1}^{T} u_{\mathrm{def}}^* - \E\!\left[\sum_{t=1}^{T} u_{\mathrm{def}}(\sigma_{\mathrm{def}}^t, \beta_{\alpha_t}(\sigma_{\mathrm{def}}^t))\right].
    \label{eq:regret_def}
\end{equation}
Equivalently, the \textbf{average regret} is $\mathrm{Regret}(T) / T$.

\subsubsection{Objective}
Our goal is to design an online algorithm for the defender such that:
\begin{enumerate}
    \item \textbf{No-regret behaviour}: The average regret vanishes as $T \to \infty$:
    \begin{equation}
        \frac{\mathrm{Regret}(T)}{T} \to 0 \quad \text{as } T \to \infty.
        \label{eq:no_regret}
    \end{equation}
    More specifically, under the assumption that $|K^*| \leq \Kmax$ and that each type $\alpha \in K^*$ has a non-zero probability of being revealed under exploration, we aim to show empirically that:
    \begin{equation}
        \mathrm{Regret}(T) = o(T),
        \label{eq:sublinear}
    \end{equation}
    i.e., the cumulative regret is sublinear in $T$, and polynomial in the game parameters $(n, k, |\mathcal{I}|)$.

    \item \textbf{Catalogue convergence}: The catalogue $C_t$ converges to $K^*$ in finite time with high probability:
    \begin{equation}
        \Pr\!\left[C_t = K^* \text{ for all } t \geq t_0\right] \geq 1 - \delta,
        \label{eq:catalogue_convergence}
    \end{equation}
    for some $t_0$ that depends on $\Kmax$, $\varepsilon_t$, and the feedback regime.

    \item \textbf{Strategy convergence}: Once $C_t = K^*$, the defender's strategy converges to $\sigma^*(K^*)$---the Nash equilibrium of the true game---as additional CFR+ iterations are performed on the stabilised game.
\end{enumerate}

%-----------------------------------------------------------------------
\subsection{Comparison with the Stackelberg Online Learning Setting}
\label{sec:comparison}
%-----------------------------------------------------------------------

Our formulation differs from Balcan et al.~\cite{Balcan2015} in several key respects:
\begin{enumerate}
    \item \textbf{Solution concept}: We seek a Nash equilibrium (simultaneous move) rather than a Stackelberg equilibrium (sequential commitment). The defender does not commit to a strategy that the attacker observes; instead, both sides act simultaneously.
    \item \textbf{Game structure}: Our game has sequential stages with resource depletion (an NFGSS with MDP-structured strategy spaces), whereas~\cite{Balcan2015} considers one-shot security games with coverage probabilities over targets.
    \item \textbf{Unknown type set}: In \cite{Balcan2015}, the set $K$ is known and the sequence of attackers is adversarial. In our setting, $K^*$ itself is unknown (with $|K^*| \leq \Kmax$), and the defender must \emph{discover} the types through interaction, not merely learn their sequence. This introduces a fundamentally different exploration requirement.
    \item \textbf{Equilibrium computation}: We use CFR+ as the inner-loop solver (re-invoked when the catalogue changes), rather than the coverage-probability-based online learning algorithms of~\cite{Balcan2015}.
\end{enumerate}

Despite these differences, the high-level explore--exploit structure---and the aspiration for sublinear regret---carries over. The regret decomposition in~\cite{Balcan2015} inspires our approach: regret accrues from (a)~exploration episodes, (b)~playing the equilibrium of an incomplete game $\Gamma(C_t) \neq \Gamma(K^*)$, and (c)~the approximation error of finite CFR+ iterations. We aim to show empirically that terms (a) and (b) are sublinear and term~(c) can be driven below any $\epsilon > 0$.

%%%%%%%%%%%%%%%%%%%%%%%%%%%%%%%%%%%%%%%%%%%%%%%%%%%%%%%%%%%%%%%%%%%%%%%%
%%% Placeholder sections for the rest of the report
%%%%%%%%%%%%%%%%%%%%%%%%%%%%%%%%%%%%%%%%%%%%%%%%%%%%%%%%%%%%%%%%%%%%%%%%

% \section{Methodology}
% \label{sec:methodology}
% TODO: Describe CFR+ implementation on NFGSS, explore--exploit algorithm details, parametric variant.

% \section{Experimental Setup}
% \label{sec:experiments}
% TODO: Resource configurations, feedback regimes, baselines (Stackelberg, oracle CFR).

% \section{Results}
% \label{sec:results}
% TODO: Convergence plots, regret curves, qualitative strategy analysis.

% \section{Conclusion}
% \label{sec:conclusion}
% TODO: Summary, limitations, future work.

%%%%%%%%%%%%%%%%%%%%%%%%%%%%%%%%%%%%%%%%%%%%%%%%%%%%%%%%%%%%%%%%%%%%%%%%

\balance

\bibliographystyle{ACM-Reference-Format}
\bibliography{references}

%%%%%%%%%%%%%%%%%%%%%%%%%%%%%%%%%%%%%%%%%%%%%%%%%%%%%%%%%%%%%%%%%%%%%%%%
\end{document}
